\documentclass[12pt]{article}
\usepackage[utf8]{inputenc}
\usepackage[T1]{fontenc}
\usepackage{amsmath,amssymb,amsfonts}
\usepackage{amsthm}
\begin{document}

APPENDIX : COHOMOLOGIE A SUPPORT PROPRE ET CONSTRUCTION DU FONCTEUR
par P. DELIGNE

(1) Verdier a montré que dans le cadre topologique, le formalisme de la dualité de Poincaré me ramenait à des probl~mes locaux en haut (voir [I]). Pour transpomer ma construction au cadre schématique, il faut disposer d'une théorie de la cohomologie "à support propre" pour les faisceaux cohérents.

Sauf mention explicite du contraire, tous les préschémas considèrent sont noethériens et les faisceaux quasi-cohérents.

\section{I. Le sorite des pro-objets}
Proposition I. Soit C une U-catégorie (2) oú existent limites inductives finies. Soit h un foncteur $C \to (\text{ens})$. Les conditions suivantes sont équivalentes :
(i) h est limite inductive, solon un petit (3) ensemble ordonné filtrant, de foncteurs représentables.
(ii) h est limite inductive, selon une petite catégorie filtrante, de foncteurs représentables.
(iii) h transforme lim finies en lim un petit ensemble d'obets telique tout élément d'un factorise par l'un d'eux.

Les implications viales et $(i) \Rightarrow (ii) \Rightarrow (iii)$ sont tristandard. Reste à prouver que (ii) $\Rightarrow$ (i).
Si $\mathcal{L},\mathcal{M}$ sont deux catégories filtrantes, un foncteur $F: \mathcal{L} \to \mathcal{N}$ est dit cofinal si $\varinjlim_{\mathcal{L}} F(L)$ est le foncteur final. Pour tout $G: \mathcal{L} \to \mathcal{N}$ on a $\varinjlim G = \varinjlim F$.

Pour toute petite catégorie filtrante, il existe un foncteur cofinal d'un petit ensemble ordonnant dans $\checkmark$.
La première projection $L \times N \to L$ est un foncteur cofinal (N muni de l'ordre naturel); ceci permet de se

(3) "petit" = " $\in U$".
ramener a supposer que $\mathrm{Ob}\,\alpha$, préordonné par $\mathrm{Hom}(x,Y) \neq\emptyset$, n'a pas de plus grand élément.
On ordonne par inclusion l'ensemble E des sous-catégories finies de $\ell$ ayant un seul objet final (fini signifie d'ensemble de flèches fini).
On définit un foncteur de E dans $\alpha$ en associant à chacune de oes cat@gories son objet final.
Sous les hypotheses faites, E est filtrant et ce foncteur cofinal.

Les foncteurs v@rifiant les conditions 6quivalentes de la proposition I sont les Ind-objets de C une sous-catégorie pleine de $\mathrm{Hom}(C,(\mathrm{Ens}))$, stable par limite inductive.
$C \to \mathrm{Ind}\,C$ est un foncteur.
Si les limites inductives filtrantes existent dans C, pour tout $h\in\mathrm{Ind}$, le foncteur $X \mapsto \mathrm{Hom}(h,h_X)$ est corepresentable, d'où un foncteur, noté $\varprojlim$, de $\mathrm{Ind}\,C$ dans C.
En particulier, on a toujours un foncteur $\mathrm{Ind}\to\mathrm{Ind}\,C$, et tout foncteur $C\to\mathrm{Ind}$ se prolonge en un foncteur $\mathrm{Ind}\,C\to\mathrm{Ind}\,C$.
Dans Ind C, les limites inductives sont exactes, et seront notées $\varinjlim$, identifiant C à une sous-catégorie pleine; en particulier, si $(x_i)_{i\in I}$ est un syst~me inductif filtrant dans C
\[
\varinjlim_{i} x_{i}
\]

Ce qui pr@c~de, sauf la prop. I (iii) et l'assertion pr@c@dente reste vrai en utilisant partout des limites de suites.
On définit par dualiti@ la cat@gorie pro C des pro-objets (sous-cat@gorie pleine de $\mathrm{Hom}(C,(\mathrm{Ens}))$

Proposition 2. Soit $X$ un préschéma quasi-compact quasi-séparé. La catégorie des faisceaux quasi-coh@rents sur $X$ est équivalente à la catégorie des Ind-objets de la catégorie des faisceaux de présentation finie sur $X$
la flèche est "lin $F_i \mapsto \varinjlim_i$".
si $F$ est de p.f., pour tout syst~me inductif filtrant, $\mathrm{Hom}(F,\varinjlim g_i)$ (par un recollement fini qui commute $\varinjlim$, on se ramène au cas affine).
Le foncteur pr@c@dent est donc pleinement fiddle. Son image est stable par limites inductives et sommes finies.
Ii suffit de prouver que pour tout $x \in X$, tout faisceau sur $X$, il existe $F$ de p.f. et une flèche de $F$ dans $G$ d@ns l'image.
sera limite d'images de faisceaux de p.f. et chacune d'elles quotient d'un faisceau de type fini.

Soit sur un voisimage quasi-compact $U$ une flèche $f: \mathscr{C} \to \mathscr{F}$, $\tau$ de p.f., s dans l'image.
Ii suffit de prolonger $\tau$ et $f$ sur $X$, ce qui revient k effectuer un prolongement de $U$

Lemme. Soit $\mathscr{C}$ un faisceau de p.f. sur un ouvert quasi-compact $U$, $\mathscr{F}$ un faisceau sur $X$, et $f$ une flèche de $\mathscr{C}$ dans $F$.
Il est possible d prolonger $\tau$ ot $f$ sur $X$, définis sur $X$.

Soit $\iota:U\hookrightarrow X$ l'inclusion, et $\sigma_1$ le produit fibré de $\iota_*\mathscr{C}$ et $\mathscr{F}$.
$\mathscr{C}_1$ prolonge $\mathscr{C}$, et s'envoie dans $\mathscr{F}$.
Si on représente comme limite de ses sous-faisceaux de type fini, comme $U$ est quasi-compact et $\mathscr{C}$ de p.f., l'un d'eux prolonge $\tau$.
Cor I. Pour qu'un foncteur contravariant additif soit représentable, il faut et suffit Qu'il soit exact à gauche et transforme lim filtrante en lim filtrante.
Cor 2. Tout faisceau de p.f. se prolonge en un faisceau de p.f.

\end{document}